\section{Introduction}
\label{Introduction}

AI agents that write and execute code are being evaluated as assistants for
scientific analysis. DiscoveryBench examines data-driven discovery workflows,
while ScienceAgentBench evaluates programs for tasks drawn from scientific
publications~\cite{majumder2025discoverybench,chen2025scienceagentbench}. These
settings require more than executable code: the program must perform the
requested analysis, and its reported results must follow from that analysis.
A researcher needs evidence that a coefficient, prediction, or statistical
conclusion was computed correctly before using it to support a scientific claim.
Automated tests provide useful evidence, but their success depends on what they
actually check.

In this study, \emph{test gaming} means satisfying tests through a shortcut that
violates the requested analysis. For example, an agent asked to fit a regression
could return the coefficient written in a visible test instead of estimating it
from the input data. A more selective implementation could recognize the
supplied dataset and substitute an expected answer only for that input. Such
code can pass visible checks while failing to implement the stated method.
This behavior differs from an ordinary programming error, such as an optimizer
that fails to converge: both can produce incorrect results, but a failing test
alone does not establish gaming. Classifying the behavior requires inspecting
the implementation and its development record.

Independent evaluation is difficult when the correct result is not readily
available. Barr et al. describe this challenge as the \emph{test oracle problem}
and survey approaches including specifications and metamorphic
testing~\cite{barr2015oracle}. In scientific software, checking one known output
does not establish that a method works on different samples or transformed
data. Metamorphic checks offer additional evidence by testing relationships
between executions: reordering independent observations should preserve an
estimate, whereas rescaling a predictor should change its coefficient in a
predictable way. These relationships also have limits. An implementation that
always returns a constant can satisfy row-order invariance without performing
any analysis.

Prior work establishes that shortcuts are a practical evaluation concern.
ImpossibleBench introduces conflicts between task specifications and tests,
making successful test exploitation identifiable, and studies behaviors
including test modification and special-casing~\cite{zhong2025impossiblebench}.
ScienceAgentBench also reports agents reading test labels instead of training
the requested model~\cite{chen2025scienceagentbench}. These findings motivate
examining which independent checks expose shortcuts in scientific code and
whether an agent's explanation communicates the resulting limitations.

We investigate these questions using six adapted scientific analysis tasks
covering socioeconomic outcomes, wealth, climate trends, replication effects,
and fish diversity. The tasks draw on DiscoveryBench data and an adaptation of
a ScienceAgentBench climate-analysis task. Our current dataset contains 190
valid agent runs. Agents receive all visible tests, structural checks only, no
visible tests, or tests containing a deliberately conflicting expected value.
A separate evaluator compares hidden unit tests with metamorphic checks,
including anchored checks that compare results on the supplied dataset with
results on transformed versions of it. Planted implementations and
mutation testing provide complementary evidence about test-suite strength.
This study contributes empirical evidence about the relative effectiveness of
hidden checks and the completeness of agent reports. Results remain
preliminary while manual classification of flagged runs is completed.

The study addresses the following research questions. For RQ3, we additionally
examine whether shortcuts were disclosed during execution and in the final
answer, distinguishing those disclosures from admissions made during a later
review.

\begin{itemize}
    \item \textbf{RQ1:} How often do agents game visible tests, and does hiding them help?
    \item \textbf{RQ2:} Which hidden checks best detect gamed code?
    \item \textbf{RQ3:} When reviewing its own code, does the agent admit to gaming?
\end{itemize}
